\section*{الفصل \ref{ch:b1:crystals-metals} --- البلورات (1): البلورة المثالية والفلزات}

\begin{solution}{exo:b1:crystals-metals:1}
المكعب البسيط: $8 \times \frac18 = 1$ ذرة، والتناسق 6. الممركز الجسم: $8 \times
\frac18 + 1 = 2$، والتناسق 8. الممركز الأوجه: $8 \times \frac18 + 6 \times
\frac12 = 4$، والتناسق 12.
\end{solution}

\begin{solution}{exo:b1:crystals-metals:2}
تتماس الذرات على امتداد الحرف: $a = 2R$، $Z = 1$، $C = \frac43\pi R^3/(2R)^3
= \pi/6 \approx 0.52$، وهو أقل من الممركز الجسم (0.68) والممركز الأوجه (0.74).
\end{solution}

\begin{solution}{exo:b1:crystals-metals:3}
$R = a\sqrt2/4 = 405.0 \times 0.3536 = \qty{143.2}{pm}$؛
$\rho = 4 \times 26.98/(6.022 \times 10^{23} \times (4.050 \times
10^{-8})^3) = \qty{2.70}{g/cm^3}$.
\end{solution}

\begin{solution}{exo:b1:crystals-metals:4}
$R = a\sqrt3/4 = \qty{185.8}{pm}$؛ $\rho = 2 \times 22.99/(6.022 \times
10^{23} \times (4.291 \times 10^{-8})^3) = \qty{0.966}{g/cm^3}$، وهو
متفق مع القيمة المقيسة \qty{0.97}{g/cm^3}.
\end{solution}

\begin{solution}{exo:b1:crystals-metals:5}
$a^3 = 4M/(\rho N_A) = 4 \times 107.87/(10.50 \times 6.022 \times 10^{23}) =
\qty{6.824e-23}{cm^3}$، فيكون $a = \qty{4.086e-8}{cm} = \qty{408.6}{pm}$
و$R = a\sqrt2/4 = \qty{144.5}{pm}$.
\end{solution}

\begin{solution}{exo:b1:crystals-metals:6}
البرهان هو برهان \cref{prop:b1:crystals-metals:hcp}. المغنيسيوم:
$c/a = 521.0/320.9 = 1.624$. والموشور الذي حجمه
$\frac{3\sqrt3}{2}a^2c$ يحوي 6 ذرات، فيكون $\rho = 6M/(N_A \cdot
\frac{3\sqrt3}{2}a^2c) = 6 \times 24.31/(6.022 \times 10^{23} \times
2.598 \times (3.209 \times 10^{-8})^2 \times 5.210 \times 10^{-8}) =
\qty{1.74}{g/cm^3}$.
\end{solution}

\begin{solution}{exo:b1:crystals-metals:7}
الذرات: $(0,0,0)$، $(\frac a2,\frac a2,0)$، $(\frac a2,0,\frac a2)$،
$(0,\frac a2,\frac a2)$. \omterm{def:b1:crystals-metals:site}{المواقع الثمانية الأوجه}: $(\frac a2,\frac a2,\frac a2)$،
$(\frac a2,0,0)$، $(0,\frac a2,0)$، $(0,0,\frac a2)$. \omterm{def:b1:crystals-metals:site}{المواقع الرباعية الأوجه}: النقاط
الثماني التي كل إحداثية منها $\frac a4$ أو $\frac{3a}{4}$. أي
\omterm{def:b1:crystals-metals:site}{موقع ثماني الأوجه} واحد وموقعان رباعيا الأوجه لكل ذرة.
\end{solution}

\begin{solution}{exo:b1:crystals-metals:8}
$R(\ce{Cu}) = 361.5\sqrt2/4 = \qty{127.8}{pm}$، $R(\ce{Ni}) = 352.4\sqrt2/4
= \qty{124.6}{pm}$: نصفا قطرين في حدود 3\,\% والبنية نفسها، فتستطيع كل
ذرة أن تحل محل الأخرى على \omterm{def:b1:crystals-metals:lattice}{الشبكة}. أما الهيدروجين، الأصغر بكثير، فلا
يستطيع إلا أن يستقر في \omterm{def:b1:crystals-metals:site}{المواقع البينية}: فهو يكوّن \omterm{def:b1:crystals-metals:alloy}{سبيكة بينية}.
\end{solution}

\begin{solution}{exo:b1:crystals-metals:9}
$R = 316.5\sqrt3/4 = \qty{137.0}{pm}$؛ $\rho = 2 \times 183.84/(6.022 \times
10^{23} \times (3.165 \times 10^{-8})^3) = \qty{19.26}{g/cm^3}$؛ وعدد الذرات في كل
\unit{cm^3}: $2/a^3 = 2/(3.165 \times 10^{-8})^3 = \num{6.31e22}$.
\end{solution}

\begin{solution}{exo:b1:crystals-metals:10}
يبعد مركز الوجه $(\frac a2,\frac a2,0)$ مسافة $a/2$ عن مركزي جسمي
الخليتين اللتين فوقه وتحته: $r_O = \frac a2 - R = \frac a2 - \frac{a\sqrt3}{4}
= \frac{2-\sqrt3}{4}a$. والنقطة $(\frac a2,\frac a4,0)$ على مسافة
$\sqrt{a^2/4 + a^2/16} = \frac{a\sqrt5}{4}$ من الرأسين $(0,0,0)$
و$(a,0,0)$ ومن مركزي الجسمين $(\frac a2,\frac a2,\pm\frac a2)$:
$r_T = \frac{\sqrt5 - \sqrt3}{4}a$. ومع $a = 4R/\sqrt3$:
$r_O = \frac{2-\sqrt3}{\sqrt3}R = 0.155\,R$
و$r_T = \frac{\sqrt5-\sqrt3}{\sqrt3}R = 0.291\,R$. \omterm{def:b1:crystals-metals:site}{فالموقع الرباعي الأوجه} هو
الأكبر في \omterm{def:b1:crystals-metals:packings}{المكعب الممركز الجسم}.
\end{solution}

\begin{solution}{exo:b1:crystals-metals:11}
$C = Z\cdot\frac43\pi R^3/a^3$ و$Z/a^3 = \rho N_A/M$، فيكون $C =
\frac{4\pi R^3\rho N_A}{3M} = \frac{\pi\rho N_A}{6M}(2R)^3$. النحاس:
$\pi \times 8.93 \times 6.022 \times 10^{23}/(6 \times 63.55) \times (2.556
\times 10^{-8})^3 = 0.740$.
\end{solution}

\begin{solution}{exo:b1:crystals-metals:12}
$a \approx (407.8 + 408.6)/2 = \qty{408.2}{pm}$؛ والكتلة المولية المتوسطة $(196.97 +
107.87)/2 = \qty{152.42}{g/mol}$؛ $\rho = 4 \times 152.42/(6.022 \times
10^{23} \times (4.082 \times 10^{-8})^3) = \qty{14.9}{g/cm^3}$.
\end{solution}

\begin{solution}{pb:b1:crystals-metals:1}
\textbf{1.} $Z = 2$؛ والتناسق 8.
\textbf{2.} على امتداد القطر الرئيسي: $a\sqrt3 = 4R$، $R = 286.65 \times
\sqrt3/4 = \qty{124.1}{pm}$.
\textbf{3.} $\rho = 2 \times 55.845/(6.022 \times 10^{23} \times (2.8665
\times 10^{-8})^3) = \qty{7.87}{g/cm^3}$.
\textbf{4.} $2/a^3 = 2/(2.8665 \times 10^{-8})^3 = \num{8.49e22}$ ذرة.
\textbf{5.} $C = \pi\sqrt3/8 = 0.680$.
\textbf{6.} تتفق الكتلتان الحجميتان المحسوبة والمقيسة في ثلاثة أرقام.
\textbf{7.} $\alpha$: $R = 289.8\sqrt3/4 = \qty{125.5}{pm}$؛ $\gamma$:
$R = 363.9\sqrt2/4 = \qty{128.7}{pm}$.
\textbf{8.} $\alpha$: $a^3/2 = \qty{12.17e6}{pm^3}$ لكل ذرة؛ $\gamma$:
$a^3/4 = \qty{12.05e6}{pm^3}$ لكل ذرة.
\textbf{9.} الطور $\gamma$ أكثف: $\rho_\alpha = 55.845/(6.022 \times 10^{23}
\times 1.217 \times 10^{-23}) = \qty{7.62}{g/cm^3}$، $\rho_\gamma =
\qty{7.70}{g/cm^3}$.
\textbf{10.} $(12.05 - 12.17)/12.17 = -1.0\,\%$: ينكمش القضيب بنحو
1\,\% حين يُسخَّن عبر التحول.
\textbf{11.} نعم: الممركز الأوجه أكثر اكتنازًا (0.74 مقابل 0.68)، وهذا يرجح على
نصف القطر الذري الأكبر: $(0.680/0.740) \times (128.7/125.5)^3 = 0.99$.
\textbf{12.} $r_C = a\sqrt3/8 = 356.68 \times 0.2165 = \qty{77.2}{pm}$.
\textbf{13.} $r_O = 0.0670 \times 289.8 = \qty{19.4}{pm}$.
\textbf{14.} $r_T = 0.1260 \times 289.8 = \qty{36.5}{pm}$.
\textbf{15.} الكربون (\qty{77}{pm}) ضعف أكبر موقع حجمًا:
فيجب أن تُبعَد ذرات الحديد المجاورة.
\textbf{16.} يكلّف التشويه طاقة كبيرة، فلا يدخل الحديد $\alpha$ إلا
عدد قليل جدًا من ذرات الكربون.
\textbf{17.} 4 \omterm{def:b1:crystals-metals:site}{مواقع ثمانية الأوجه} في خلية من 4 ذرات: موقع واحد لكل ذرة حديد.
\textbf{18.} $r_T = (\sqrt{3/2} - 1) \times 128.7 = \qty{28.9}{pm}$.
\textbf{19.} كربون واحد لكل حديد: \ce{FeC}، $w(\ce{C}) = 12.011/(55.845 +
12.011) = 17.7\,\%$.
\textbf{20.} كربون واحد لكل عشر ذرات حديد: $w = 12.011/(558.45 + 12.011) =
2.1\,\%$.
\textbf{21.} \omterm{def:b1:crystals-metals:site}{المواقع الثمانية الأوجه} في الحديد $\gamma$ أكبر بكثير من
أي موقع في الحديد $\alpha$: فالحديد الساخن يذيب الكربون، ثم يحبسه
السقي.
\textbf{22.} $r_O = (\sqrt2 - 1) \times 128.7 = \textbf{\qty{53}{pm}}$، وهو
أكبر فجوة في الحديد ومع ذلك أصغر من ذرة الكربون (\qty{77}{pm}):
فالكربون يذوب في الحديد الساخن، لكنه يمطّه.
\end{solution}
